% Standalone paper-level proof note. No new Lean theorem is claimed here.
\documentclass[11pt]{article}
\usepackage[T1]{fontenc}
\usepackage{lmodern,amsmath,amssymb,amsthm,geometry,booktabs,xurl}
\usepackage[hidelinks]{hyperref}
\geometry{margin=1in}
\setlength{\emergencystretch}{2em}
\newtheorem{theorem}{Theorem}
\newtheorem{lemma}[theorem]{Lemma}
\newtheorem{proposition}[theorem]{Proposition}
\newcommand{\PA}{\mathrm{PA}}
\newcommand{\Bit}{\operatorname{Bit}}
\newcommand{\Prf}{\operatorname{Prf}}
\newcommand{\Con}{\operatorname{Con}}
\newcommand{\I}{\mathsf I}
\newcommand{\M}{\mathsf M}
\newcommand{\U}{\mathsf U}
\newcommand{\E}{\mathsf E}
\newcommand{\N}{\mathsf N}
\newcommand{\tr}{\tau}
\title{Internal soundness of the PA196 proof search:\\a direct syntactic route into Peano arithmetic}
\author{Research note}
\date{13 November 2026}
\begin{document}
\maketitle

\begin{abstract}
The published PA196 development proves in PA that literal nonhalting
implies $\Con(\PA)$, but its reverse soundness argument is external.
This note gives a paper-level construction for the missing direction:
translate the machine's numeric formulas directly into PA, compile each
axiom instance into a PA proof, and retain proof witnesses beside the
machine's formula registers.  The crucial nonlogical template is finite
bit comprehension proved by one full-induction instance.  No truth
predicate, interpretation through the older M6 decoder, or equality of
old and recoded binder labels is used.  We spell out the positive-length
simulation argument needed to extract a source computation from a literal
halting trace.  This is a mathematical proof construction, not a new
kernel-checked Lean certificate.
\end{abstract}

\section{Exact scope and upstream status}

All source references in this note refer to James Chang's
\href{https://github.com/jamesyc/BusyBeaver/tree/407cdc356f6b586635e71ac92be394255f4bfea4}{PA196 revision
\nolinkurl{407cdc356f6b586635e71ac92be394255f4bfea4}}.
Its theorem \texttt{pa\_proves\_nonhalting\_consistency} proves
$\PA\vdash\neg H\to\Con(\PA)$, where $H$ is the exact arithmetic
blank-tape halting sentence.  The theorem
\texttt{machine\_nonhalting} is an external theorem in Lean's metatheory.
Neither declaration, by itself, supplies $\PA\vdash H\to\neg\Con(\PA)$.
The difference matters for nonstandard proof codes and nonstandard
computation lengths.

We use the ordinary, effective first-order proof predicate for PA.  A
conservative definitional extension by the elementary arithmetic
operations used below is harmless: each auxiliary proof can be
translated back into the original language by a primitive recursive
proof translation whose correctness is provable in PA.  The chosen ordinary
PA proof calculus must likewise be calibrated to the upstream PA predicate:
standard Hilbert/sequent presentations admit primitive recursive mutual
proof translations, provably correct in PA.  The consistency formula in
the conclusion is the upstream one after this fixed calibration.  In particular,
$2^n$ and binary-digit tests can be treated as abbreviations for their
standard arithmetic graphs; no exponentiation axiom is added to PA.

\section{Translation without recoding variable names}

Let $p(a,b)=\binom{a+b+1}{2}+b$ and use exactly
\[
 \I(a,b)=4p(a,b),\quad \M(a,b)=4p(a,b)+1,\quad
 \U(a,B)=4p(B,a)+2,\quad \E(a,b)=4p(a,b)+3.
\]
The decoder gives code $0$ the formula $\top$.  For positive code
$n=4p(u,v)+r$, its constructor is implication, membership, universal
quantification, or equality according as $r=0,1,2,3$; the universal case
binds variable \emph{$v$} and has body code $u$.

Define $\tr(n)$ directly as an ordinary arithmetic formula:
\[
\begin{array}{rcl}
 \tr(0)&=&0=0,\\
 \tr(4p(u,v))&=&\tr(u)\to\tr(v)\qquad(u+v>0),\\
 \tr(\M(u,v))&=&\Bit(x_u,x_v),\\
 \tr(\U(v,u))&=&\forall x_v\,\tr(u),\\
 \tr(\E(u,v))&=&x_u=x_v.
\end{array}
\]
Here $\Bit(A,B)$ says that the $B$th binary digit of $A$ is $1$.
The numeral indexing a variable is retained unchanged.  In particular,
if the source gate binds $x_c$, the translated gate binds $x_c$ even
though the G\"odel code of $\tr(c)$ will generally not be $c$.

\begin{lemma}[Syntactic bounds]
The translation is primitive recursive.  PA proves, uniformly in $n$,
that every source-variable label occurring in $\tr(n)$ is strictly
less than $n$.  Consequently $x_n$ does not occur in $\tr(n)$.
\end{lemma}
\begin{proof}
For $n>0$, both components of the inverse Cantor pairing of
$\lfloor n/4\rfloor$ are at most $\lfloor n/4\rfloor<n$.
Course-of-values recursion on $n$ therefore defines $\tr$; it can
instead be implemented by a table of the values at $0,\ldots,n$.
The same induction proves the bound on the labels at every atomic
leaf and at every quantifier.  The argument is an induction on finite
\emph{syntax codes}; it is not an induction defining arithmetic truth.
A literal expansion of $\Bit$ may introduce bound auxiliary variables.
Choose these from a fresh block above all source labels, and alpha-rename
them when embedding a translated subformula.  The assertion concerns
source labels; the usual free-variable bound, which is the one needed
for capture avoidance, is unaffected by these bound auxiliaries.
\end{proof}

There are two small but essential boundary cases.
First, $\I(0,0)=0$, so literal compositionality for implication fails only
at this pair.  Uniformly in $a,b$, generate a PA proof of
$\tr(\I(a,b))\leftrightarrow(\tr(a)\to\tr(b))$: outside the
exception the formulas coincide; at $(0,0)$ use the fixed tautology
$\top\leftrightarrow(\top\to\top)$.  Formula congruence transports
this equivalence through the fixed shapes of the axiom heads, so a nested
zero collapse (for example in head 6 or $\rho_8$) is covered too.
A proof translation for modus ponens branches on this
case and returns the fixed PA proof of $0=0$.
Second, code $1=\M(0,0)$ is not syntactically the falsum symbol; its
translation is $\chi:=\Bit(x_0,x_0)$.  PA proves
\[
 \forall z\,\neg\Bit(z,z)
\]
from $2^z>z$ and the elementary digit bound.  Thus a PA proof of
$\tr(1)$ can be converted to a PA refutation.  Everywhere the source
uses $\N B:=\I(B,1)$, source conjunction, or source existential
quantification, use the fixed ordinary logical derivations justified
by $\forall z\neg\Bit(z,z)$.  They remain valid if a displayed
quantifier binds $x_0$; one must not treat $\chi$ as a syntactically
closed atom.

\section{A uniform PA proof of the guarded gate}

\begin{lemma}[Uniform finite bit comprehension]\label{lem:comprehension}
For each arithmetic formula $R(z,\vec a)$, PA proves
\[
 \forall n\ \exists s<2^n\ \forall z<n\,
       \bigl(\Bit(s,z)\leftrightarrow R(z,\vec a)\bigr).
\]
Moreover, there is a primitive recursive procedure that, given the
syntax code of $R$ and a fresh-variable list, outputs a PA proof of
this sentence; PA verifies the correctness of that procedure on
all formula codes.
\end{lemma}
\begin{proof}
Use the induction formula
\[
 Q(n):=\exists s<2^n\ \forall z<n\,
       \bigl(\Bit(s,z)\leftrightarrow R(z,\vec a)\bigr).
\]
For $n=0$ take $s=0$.  Given a witness $s<2^n$ for $Q(n)$, apply
excluded middle to $R(n,\vec a)$.  In its false case retain $s$;
in its true case use $s+2^n$.  Elementary arithmetic proves that
addition of $2^n$ changes just the $n$th bit of a number below $2^n$,
that lower bits are unchanged, and that the new number is below
$2^{n+1}$.  These are fixed PA theorems.  Hence the step is a fixed
logical proof template with $R$ and its capture-avoiding instances
substituted in it.  Append exactly the induction axiom for $Q$ and
modus ponens.

To see the claimed uniformity, implement the generator by term and
formula constructors, freshening, capture-avoiding substitution,
proof concatenation, and line-number shifts.  Each is primitive
recursive.  The output contains one induction instance for the
input formula, not a decision whether that formula is true.
PA checks the output by inspecting the local proof rules and the
syntactic induction-axiom test.  Induction on the input syntax
verifies substitution and freshening.  Consequently this is a
uniform proof \emph{compiler}, not a uniform satisfaction predicate.
\end{proof}

Write $\varphi=\tr(c)$, and assume the machine's guard $c\ne a,b$.
The syntactic-bounds lemma ensures $x_c$ is fresh in $\varphi$.
Up to the fixed translations of the derived connectives, pointed
separation is
\begin{align*}
 \exists x_c\,(S\wedge P),\qquad
 P&:=\varphi\to\Bit(x_c,x_b),\\
 S&:=\forall x_b\bigl[
    (\Bit(x_a,x_b)\to(\varphi\to\Bit(x_c,x_b)))
    \wedge(\Bit(x_c,x_b)\to\varphi)\bigr].
\end{align*}
To compile a proof, introduce fresh variables for a copy $z$ of $x_b$
and the selected-set code $s$, and apply Lemma~\ref{lem:comprehension}
with
\[
 n=x_a+x_b+1,\qquad
 R(z):=(\Bit(x_a,z)\vee z=x_b)\wedge
          \varphi[z/x_b].
\]
The substitution replaces only \emph{free} occurrences of $x_b$;
all auxiliary binders are fresh.  The variables $x_a,x_b$ in the
first parenthesis here are the outside parameters.
Every set bit of $x_a$ has index below $n$, and $x_b<n$.  The resulting
$s$ therefore contains every member of $x_a$ satisfying the property,
contains no member failing it, and contains the distinguished point
$x_b$ whenever that point satisfies it.  Taking $x_c=s$ proves
$S\wedge P$.

If $a=b$, the occurrence of $x_a$ inside $S$ is also bound by its
$\forall x_b$.  Then the first antecedent becomes $\Bit(x_b,x_b)$
and is false.  The same selected set still proves the second
conjunct and the point clause.  A proof generator branches on
$a=b$ and uses the fixed irreflexivity proof for this case; it must
not silently regard $x_a$ as a fixed parameter inside $S$.

Under $c\ne a,b$, ordinary logical vacuity supplies proofs of
\[
 \Delta:=x_a=x_b\to\forall x_c(x_a=x_b),\qquad
 V:=\Bit(x_a,x_b)\to\forall x_c\Bit(x_a,x_b).
\]
Let $Q:=\forall x_c(S\to\neg P)$.  The pointed-separation proof gives
$\neg Q$.  Propositional logic then gives
\[
 F(a,b,c)=\neg\bigl(\Delta\to(V\to Q)\bigr).
\]
This is precisely the translated PA196 gate, modulo the fixed
code-$1$ connective translations.  If the guard fails, the machine
returns code $0$, whose translation has a fixed proof.

\section{Compiling proofs and machine commands}

\begin{theorem}[Uniform syntactic soundness]
There is a primitive recursive proof compiler $J$ such that PA proves
\[
 \operatorname{NativeProof}(d,n)\ \longrightarrow\
 \Prf_{\PA}(J(d),\ulcorner\tr(n)\urcorner).
\]
In particular, PA proves that a native proof of code $1$ gives a PA
refutation.  The same conclusion holds for the exact fuelled PA196
checker by induction on its fuel table.
\end{theorem}
\begin{proof}
Compile each of the six heads with an ordinary first-order logical
proof template.  Heads 3 and 4 are equality substitution for the
formula $\Bit$; head 5 is the existential identity witness.  Head 1
is $(\exists x\forall x A)\to A$, head 2 exchanges universal
quantifiers, and head 6 distributes universal quantification over
modus ponens.  Coincident labels are permitted: ordinary formula
substitution and alpha-renaming supply the correct instances.
The restricted head $\rho_8$ is a propositional tautology for all
four parameters, so its additional numerical guard is harmless to
soundness.  Compile the guarded gate by the preceding section.
Generalization appends universal generalization to a theorem proof;
modus ponens concatenates its two translated premise proofs and
adds the MP line, except in the $\I(0,0)=0$ case already discussed.
The proof-checker recursion has exactly these constructors.  Induction
on a proof tree, or on the fuel rounds of the specified table checker,
verifies each output proof by the elementary syntactic checks.
\end{proof}

There is a second route that avoids binding a separate native-proof
predicate to the machine: attach a PA-proof witness directly to
each live formula.  A store $(t,w,d,h)$ carries proofs of
$\tr(t)$ and $\tr(h)$, plus a parallel stack of proofs for the
formula stack encoded by $w$.  The zero tail has an infinite implicit
tail of code-$0$ formulas with the fixed truth proof.  On an AX
command, push the saved proof, append generalization of the proof
of $h$ at label $d$, and append the eight uniformly compiled
axiom/gate proofs.  On successful MP append the translated MP
proof; on failed MP use the truth proof.  Shift only moves proof
witnesses.  If its test $t+h=1$ succeeds, natural arithmetic gives
$t=1$ or $h=1$, and the corresponding retained proof yields a PA
refutation.

All augmented operations are primitive recursive on codes.  PA
proves their invariant by induction on the length of a coded
command execution.  This proves
\[
 \PA\vdash
 \operatorname{CommandHalt}(0,0,0,0)\to
 \exists e\,\Prf_{\PA}(e,\ulcorner\bot\urcorner).
\]
The argument does not assume native nonrefutability, standard-model
soundness, or the existence of an arithmetic truth predicate.

\section{Extracting a command history from a literal halt}

A forward simulation alone is insufficient unless its positive
progress and nonhalt endpoints are retained.  The following elementary
lemma supplies the reverse direction uniformly over PA models.

\begin{lemma}[Bounded reverse extraction]\label{lem:reverse}
Suppose two deterministic, arithmetically coded transition systems
have a PA-provable refinement relation $R$ with these properties:
\begin{enumerate}
\item A fixed finite startup reaches a target configuration related to
the initial source configuration.  It has no halt before its endpoint.
\item A related source configuration is either terminal, or has a
positive-length target segment to a configuration related to its
source successor.  The endpoint is a working configuration whenever
the successor source configuration is live, and every intermediate edge
is a genuine target edge.  A live current state is allowed to have a
terminal successor.
\item Target halt either has no outgoing edge or is absorbing under the
total target step.  Coded source traces of every specified finite length
exist, target step is functional, and every endpoint related to a live
source configuration is nonhalted.
\end{enumerate}
Then PA proves that a target halt implies a source halt.  A target
halt at $T$ steps forces a terminal source configuration within at
most $T+1$ source transitions after startup.
\end{lemma}
\begin{proof}
All statements can be written in arithmetic using coded finite
vectors.  Fix a putative target trace of length $T$ to halt.  A halt
before the fixed startup endpoint contradicts functionality and the
startup trace.  Suppose the initial source configuration and the configurations
after its first $T+1$ transitions are all nonterminal.  Inductively concatenate their target segments.
After $i$ source steps the accumulated length is at least $i$ and
the endpoint is a working configuration.  This is a PA induction on
an arithmetically expressible existence-of-coded-segment invariant.
At $i=T+1$ it gives a target trace longer than $T$.  Determinism
forces its configuration at time $T$ to be the halt of the given
trace.  If halt has no outgoing edge, its next edge is impossible.
If halt is absorbing, induction from time $T$ forces the endpoint of
the longer trace to be halted, contradicting its live endpoint.
The same proof is valid for nonstandard $T$ and coded vectors
in any PA model, because the induction is an instance of PA's
induction scheme, not an external induction on standard naturals.
\end{proof}

For PA196, the necessary strengthening of the published
\texttt{physical\_step} theorem is already visible in its proof:
its left branch is taken when the resolved instruction is literally
\texttt{halt}; call that pre-execution configuration terminal for
Lemma~\ref{lem:reverse}.  One further source step reaches the source
halted sentinel.  Its right branch is taken otherwise and produces
\texttt{Progress true} and the successor \texttt{Refines} relation.
Retain this branch condition, instead of existentially forgetting it.
The exact control labels in this part of the development are the
auxiliary halt $215$ and the literal halt $196$.  Project each
positive segment through the finite state map; the tape and segment
length are unchanged.  The active-state invariant excludes the
nineteen unused auxiliary controls.  This uses the checked row
binding and cannot identify the two halt numerals by equality.

For the source-to-command layer one must use the absorbing version of
Lemma~\ref{lem:reverse}: the coded source step satisfies
\texttt{InternalCodedSourceStep.step\_halted}, so control $0$ has a
self-loop.  It is incorrect to claim that this coded halt has no
outgoing source edge.  Positive progress is also essential: the
existing \texttt{source\_consume\_step} may take zero steps when the
residual is $0$, so it cannot directly support the counting argument.

Instead use the total global scheduler on triples $(N,L,S)$, where
$N$ is the candidate register, $L$ the residual, and $S$ is the HFS
encoding of the four-register command store.  The empty encoding
$S=0$ means halt; a live command store is a nonempty vector and can
never have that encoding.  For a live triple define
\[
 \begin{array}{c|cc}
       &L=0&L>0\\ \hline
 \mathrm{raw}&N&L-1\\
 N'&N+1&N
 \end{array}
 \qquad
 L'=\operatorname{residual}(\mathrm{raw}),\qquad
 S'=\operatorname{CommandStep}
       (\operatorname{commandAtom}(\mathrm{raw}),S).
\]
Set $F(N,L,S)=(N',L',S')$, and make triples with $S=0$ absorbing.
All functions have the explicit arithmetic graphs of
\texttt{SourceCommandDecodeInternal.lean} and
\texttt{CommandStep.lean}; thus $F$ and its finite iterations are
arithmetically representable and PA-provably total.

Relate this scheduler state to a coded source configuration by the
exact \texttt{CommandRelation} of
\texttt{InternalCodedSourceCommandRelation.lean}:
\[
 (S=0\ \wedge\ \operatorname{control}(C)=0)\quad\vee\quad
 \begin{gathered}
 \operatorname{control}(C)=26,\quad
 \text{the source has a valid clean nine-register store},\\
 r_0=N,\quad r_8=L,\quad
 [r_1,r_2,r_4,r_5]=S.
 \end{gathered}
\]
The clean condition is $r_3=r_6=r_7=0$.  For $S\ne0$ only the
second branch is possible.  When $L=\mathrm{raw}+1$, the already
proved all-model theorem
\nolinkurl{InternalCodedSourceRootCommand.root_positive_command}
gives a source segment of positive length to a configuration related
to $F(N,L,S)$.  When $L=0$, the theorem
\texttt{root\_fresh\_command} gives the same conclusion and increments
$N$.  Both theorems use
\nolinkurl{InternalCodedSourceDispatchTotal.dispatch_command}, which
handles the three command classes.  The results use no
\texttt{not\_refutable} premise.  If $S'\ne0$, the endpoint has
control $26$ and is live; if $S'=0$, the nonempty-vector fact rules
out the live branch, so its control is $0$.

The clean initial triple is
$(2,0,[0,0,0,0])$.  Suppose its coded source execution halts in $T$
steps but its initial scheduler state and all states after its first
$T+1$ scheduler iterations stay live.  Iterating
the appropriate positive root theorem gives a source trace of length
at least $T+1$ with control $26$ at its endpoint.  Determinism and
absorption at source control $0$ force the same endpoint to have
control $0$, a contradiction.  Hence some scheduler step produces
$S'=0$.  A live command state can become $0$ only in Shift's
$t+h=1$ branch, as is explicit in \texttt{CommandStep.stepExpr} and
\texttt{shift\_value}; MP and AX always return four-entry nonempty
vectors.  Retaining the finite sequence of scheduler command atoms
therefore supplies exactly the command history required by the
parallel PA-proof-witness invariant.  That invariant yields a PA
refutation.

\subsection{A shorter proof under the consistency hypothesis}

The reverse-extraction construction is useful as a reusable lemma,
but the desired implication admits a more direct proof.  Work
\emph{inside PA} and assume $\Con(\PA)$.  This is the antecedent of
the implication being proved, not an external soundness assumption.
The parallel proof-witness invariant shows that every scheduler
command remains live: a successful Shift halt would supply a PA
proof of $\tr(1)$ and hence a PA refutation, contrary to the assumed
consistency sentence.

Starting from $(2,0,[0,0,0,0])$, use the appropriate positive root
command theorem at each step.  PA induction produces, for every
$n$, a coded source trace of length at least $n$ ending at a clean
root configuration of control $26$ and carrying the proof witnesses.
Because coded source halt is absorbing, such a trace never passes
through control $0$.  For each source step, retain the nonhalt
branch of \texttt{physical\_step}; its proof supplies positive
physical progress whenever the instruction is not a halt.  An
instruction that were halt would make the next source control $0$,
which has just been excluded.  Concatenate these physical segments,
prepend startup, and project their active controls to the literal
196-state table.  For every $n$, this gives a genuine coded literal
trace of length at least $n$ with a working endpoint.

If a literal halt occurred at time $T$, choose $n=T+1$ and compare
that trace with the longer one.  Literal determinism identifies the
configuration at time $T$ as halt, while the longer trace has a
following genuine transition; the literal halt has no such
transition.  This contradiction is obtained by arithmetic induction
on coded traces and applies also to nonstandard $T$.  Consequently
\[
 \PA\vdash\Con(\PA)\to\neg H.
\]
Classical contraposition gives the desired reverse direction
$\PA\vdash H\to\neg\Con(\PA)$.  Together with the published
internal implication it yields the paper-level conclusion
\[
 \PA\vdash\neg H\leftrightarrow\Con(\PA).
\]

The concrete proof compiler and the strengthened positive simulation
lemmas have \emph{not} been entered and kernel-checked in the upstream
Lean development by this note.  An implementation should export
these declarations explicitly; it should not rename the existing
one-way theorem or cite external HF soundness as their proof.

\section{What the iPRA work can and cannot contribute}

The ConciseIPRA archive studies proof presentations equivalent to
intuitionistic primitive recursive arithmetic.  Replacing the PA196
calculus by one of them changes the consistency strength; consistency
of iPRA cannot be substituted for consistency of full PA.  The useful
transfer is architectural: explicit syntactic constructors, small
local certificate rules, retained witnesses, and elimination of
redundant axiom heads by actual proof templates.  A size comparison
of Python checkers is not a state-count bound for the binary tape
machine, and does not by itself validate a new TM construction.

\section{Membership-head deletion needs a different argument}

The two direct membership heads are
\[
 \E(a,b)\to(\M(d,a)\to\M(d,b)),\qquad
 \E(a,b)\to(\M(a,d)\to\M(b,d)).
\]
Deleting one would shorten the emitter substantially.  An ordinary
first-order semantic argument using separation suggests that
incoming-neighborhood indistinguishability might force outgoing
indistinguishability.  That is not yet a derivation in this restricted
numeric calculus.  The published reverse-rewriting template for a
property $\chi(t)$ uses the corresponding \emph{forward} rewriting
rule to specialize the separated negative property.  Applying that
same template to the missing argument position is circular.
The companion note \texttt{membership\_independence.tex} strengthens this
observation: explicit counterinterpretations show that neither head is
derivable from the remaining exact native rules, even with the full
guarded gate. A different PA interpretation of a weaker one-head calculus
is a separate question and is not excluded by those counterinterpretations.
No deletion of either head is justified by the present construction.

\end{document}
